\chapter{Sistemas ortonormales completos}

\begin{definicion}
    Sea $S\subseteq \mathcal{H}$ se dice \tbf{sistema ortonormal} si está constituido por vectores mutuamente ortogonales y de norma unitaria:
    \[ S = \{ e_i \in \mathcal{H}, \ i \in \mathcal{I} : (e_i, e_j) = \delta_{ij}\}\]
\end{definicion}


\begin{ejemplo}
    \begin{enumerate}
        \item $\ell_2 : \{e_i\}_{i=1}^{+\infty}$ es un sistema ortonormal.
        \item $L^2((0, \pi)) : \left\{ \frac{1}{\sqrt{2\pi}}, \frac{1}{\sqrt{\pi}} \cos(\pi t), \frac{1}{\sqrt{\pi}} \sin(\pi t) \right\}_{t=1, \dots, +\infty}$ es un sistema ortonormal de vectores. \\
    \end{enumerate}
\end{ejemplo}

\observacion{
    Si $S$ es un sistema ortonormal, entonces $S$ es un sistema de vectores linealmente independientes.
}

\observacion{
    Si $S$ es un sistema de vectores linealmente independientes, entonces se puede construir una sucesión ortonormal mediante el proceso de \textbf{Gram-Schmidt}. \\
}

\begin{definicion}
    Un sistema ortonormal se dice \textbf{completo} si no está contenido propiamente en ningún otro sistema ortonormal (es decir, es maximal por inclusión). \\
\end{definicion}

\begin{prop}
    Sea $\mathcal{H}$ un espacio con producto interno, entonces por el \textbf{Lema de Zorn} siempre existe un sistema ortonormal completo. 
\begin{proof}
    Sea 
    $$S = \{ e_i = \frac{x}{\|x\|}, x \neq 0 , x\in \mathcal{H}\}$$ 
    y sea $\mathcal{P}$ la familia de todos los conjuntos ortonormales que contienen a $S$ con la relación de orden dada por la inclusión. Se puede demostrar que $\mathcal{P}$ es inductiva y, por lo tanto, por el Lema de Zorn, admite un elemento maximal.
\end{proof}
\end{prop}

\begin{ejercicio}
    Sea $\mathcal{H}$ un espacio de Hilbert, entonces 
    \begin{equation*}
        S \text{ es un sistema ortonormal completo} \iff \text{ si } \forall x \in \mathcal{H}, x \perp S \text{ entonces } x = 0
    \end{equation*}
    \begin{description}
        \item[$\Longrightarrow)$] Sea $x \in \mathcal{H} : (x, s) = 0 \quad \forall s \in S$. Entonces, o bien $x=0$, o bien 
        $$\frac{1}{\|x\|}(x, s) = \left(\frac{x}{\|x\|}, s\right) = 0 \implies \frac{x}{\|x\|} \perp S$$
        Por lo tanto, $S \cup \{ \frac{x}{\|x\|} \}$ sería aún un sistema ortonormal, lo cual contradice la maximalidad de $S$ (a menos que $x=0$).
        \item[$\Longleftarrow)$] Si $x \perp S$, entonces $x = 0$, esto implica que no existe ningún sistema ortonormal $S' \supsetneq S$. Por lo tanto, $S$ es completo.
    \end{description}
\end{ejercicio}

\begin{ejercicio}
    $S$ es completo $\iff \mathcal{H} = \overline{[S]}$ (la clausura del subespacio generado por $S$)
\end{ejercicio}

\begin{prop}[Desigualdad de Bessel finita]
    Sea $\{e_i\}_{i=1}^{n}$ un sistema ortonormal en $\mathcal{H}$. Entonces, $\forall x \in \mathcal{H}$:
    \[ \sum_{i=1}^{n} (x, e_i)^2 \le \|x\|^2 \]
    Además, 
    $$x - \sum_{i=1}^{n} (x, e_i)e_i \perp e_j \quad \forall j=1, \dots, n$$
\begin{proof}
    Sea $S_n := \sum\limits_{i=1}^{n} (x, e_i)e_i$, entonces 
    $$0 \le \|x - S_n\|^2 =(x-S_n, x-S_n)= \|x\|^2 + \|S_n\|^2 - 2(x, S_n)$$
    y por otro lado tenemos que 
    \begin{align*}
        \|S_n\|^2 &= \left( \sum_i (x, e_i)e_i, \sum_j (x, e_j)e_j \right) = \sum_{i=1}^n \sum_{j=1}^n (x, e_i) (x, e_j) (e_i, e_j)\overset{(*)}{=} \sum_i (x, e_i)^2 \\
        (x, S_n) &= \left( x, \sum_i (x, e_i)e_i \right) = \sum_i (x, e_i)^2
    \end{align*}
    donde en $(*)$ se ha utilizado la ortonormalidad de los $e_i$, es decir, $(e_i, e_j) = \delta_{ij}$. Tenemos finalmente la desigualdad buscada
    \begin{equation*}
        0 \le \|x\|^2 - \sum_{i=1}^{n} (x, e_i)^2 \implies \sum_{i=1}^{n} (x, e_i)^2 \le \|x\|^2
    \end{equation*}
    Para la segunda parte, tenemos que
    \begin{equation*}
        \left( x - \sum_{i=1}^{n} (x, e_i)e_i, e_j \right) = (x, e_j) - \sum_{i=1}^{n} (x, e_i) \delta_{ij} = (x, e_j) - (x, e_j) = 0
    \end{equation*}
    lo que implica que $x - \sum\limits_{i=1}^{n} (x, e_i)e_i \perp e_j$ para todo $j=1, \dots, n$.
\end{proof}
\end{prop}

\begin{prop}
    Sea $\mathcal{H}$ un espacio de Hilbert y sea $S = \{e_i\}_{i \in \mathcal{I}}$ un sistema ortonormal. Entonces, para cada $x \in \mathcal{H}$:
    \[ S(x) := \{ i \in \mathcal{I} \mid (x, e_i) \neq 0 \} \text{ es a lo sumo numerable} \]
\begin{proof} % //TODO: leer demostración
    Sea $x \in \mathcal{H} : x \neq 0$. Para cada $n \ge 1$, definimos:
    \[ \mathcal{S}_n := \left\{ i \in \mathcal{I} \mid |(x, e_i)|^2 > \frac{\|x\|^2}{n} \right\} \]
    Entonces, $\mathcal{S}_n$ tiene a lo sumo $n-1$ elementos. De hecho, supongamos por el contrario (P.A.) que $|\mathcal{S}_n| \ge n$, entonces:
    \[ \sum_{i=1}^{n} |(x, e_i)|^2 > \sum_{i=1}^{n} \frac{\|x\|^2}{n} = \frac{n}{n} \|x\|^2 = \|x\|^2 \]
    Lo cual es un \textbf{ABSURDO} por la desigualdad de \textbf{Bessel finita}.
    \noindent
    Como $S(x) = \bigcup_{n=1}^{\infty} \mathcal{S}_n$, entonces $S(x)$ es numerable (unión numerable de conjuntos finitos). \\
\end{proof}
\end{prop}

\begin{definicion}
    Sea $S \subseteq \mathcal{H}$ un sistema ortonormal completo, dado por $S = \{e_i\}_{i \in \mathcal{I}}$, entonces para cada $x \in \mathcal{H}$, definimos la serie:
    \[ \sum_{i \in \mathcal{I}} |(x, e_i)|^2 \quad \text{con } (x, e_i) := \text{coeficientes de Fourier de } x \text{ respecto a } S \]

\end{definicion}

\observacion{
    Si demostramos que $\sum_{i \in \mathcal{I}} |(x, e_i)|^2$ converge, entonces converge absolutamente y, por lo tanto, podemos permutar los índices.
}


\begin{teo}[Desigualdad de Bessel]
    Sea $S = \{e_i\}_{i \in \mathcal{I}}$ un sistema ortonormal, entonces para cada $x \in \mathcal{H}$: 
    \[ \sum_{i \in \mathcal{I}} |(x, e_i)|^2 \le \|x\|^2 \]
\begin{proof}
    \begin{description}
        \item[$S(x)$ finito.] Ya está demsotrado.
        \item[$S(x)$ numerable.] Numeramos los elementos para los cuales $(x, e_i) \neq 0$ como $\{e_1, e_2, \dots\}$, entonces
        \[ \forall m : \sum_{i=1}^{m} |(x, e_i)|^2 \le \|x\|^2 \xrightarrow{m \to +\infty} \sum_{i=1}^{+\infty} |(x, e_i)|^2 \le \|x\|^2 \]
    \end{description}
    \noindent
    Mostremos que no depende de la elección del orden de los índices. \\

    \noindent
    Consideremos la serie $\sum_{i \in \mathcal{I}} (x, e_i)e_i$ y sea $\{e_1, \dots, e_n, \dots\}$ tal que $(x, e_i) \neq 0$. Sea 
    $$S_m := \sum_{i=1}^{m} (x, e_i)e_i$$
    entonces para $n \ge m$ se tiene:
    \[ \|S_n - S_m\|^2 = \left\| \sum_{i=m+1}^{n} (x, e_i)e_i \right\|^2 \overset{(*)}{=} \sum_{i,j=m+1}^{n} (x, e_i)\overline{(x, e_j)}(e_i, e_j) = \sum_{i=m+1}^{n} |(x, e_i)|^2 \]
    donde en $(*)$ se ha utilizado la conjugación de la segunda componente del producto interno para $\C$ (para $\R$ es el mismo número). Tenemos a continuación que la norma anterior es convergente $\forall n, m$, y por lo tanto:
    \[ \|S_n- S_m\| \xrightarrow{m, n \to +\infty} 0 \implies \{S_m\} \text{ es de Cauchy} \overset{(*)}{\implies} \exists y \in \mathcal{H} : y = \sum_{i=1}^{+\infty} (x, e_i)e_i \]
    donde en $(*)$ se ha utilizado la completitud de $\mathcal{H}$. \\

    \noindent
    Sea $\{e'_1, e'_2, \dots, e'_m, \dots\}$ una numeración diferente, entonces $\exists y' := \sum\limits_{i=1}^{+\infty} (x, e'_j)e'_j$. Mostremos que 
    $$y = y' \iff \|y - y'\| < \varepsilon \ \forall \varepsilon > 0$$
    Sea $\varepsilon > 0$, entonces 
    $$\exists \nu_\varepsilon > 0 \text{ tal que } \forall m > \nu_\varepsilon \text{ que cumple } \|S_m - y\| < \varepsilon, \ \|S'_m - y'\| < \varepsilon, \sum_{i=\nu_\varepsilon+1}^{+\infty} |(x, e_i)|^2 < \varepsilon^2$$
    donde se ha usado la definición de límite para $y,y'$, y que la serie de Bessel converge al estar acotada por $\|x\|^2$. \\

    \noindent
    Fijamos $m > \nu_\varepsilon$, entonces $\exists m' > \nu_\varepsilon : \{e_1, \dots, e_m\} \subseteq \{e'_1, \dots, e'_{m'}\}$ y por lo tanto:
    \[ \|S'_{m'} - S_m\|^2 = \left\| \sum_{\text{finita}} (x, e_i)e_i \right\|^2 =\sum_{\text{finita}} |(x, e_i)|^2\le \sum_{i=\nu_\varepsilon+1}^{+\infty} |(x, e_i)|^2 < \varepsilon^2 \]
    \[ \|y - y'\| \le \underbrace{\|y - S_m\|}_{\le \varepsilon} + \underbrace{\|S_m - S'_{m'}\|}_{\le \varepsilon} + \underbrace{\|S'_{m'} - y'\|}_{\le \varepsilon} \le 3\varepsilon \quad \forall \varepsilon > 0 \implies y = y' \]
\end{proof}
\end{teo}

\begin{definicion}
    Sea $S$ un sistema ortonormal en un espacio de Hilbert $\mathcal{H}$, entonces $\forall x \in \mathcal{H}$ está definida la suma de la serie:
    \[ \sum_{i \in \mathcal{I}} (x, e_i) e_i \quad \text{denominada \textbf{Serie de Fourier} de } x \text{ en } \mathcal{H} \]

\end{definicion}

\begin{prop}
    Sea $S$ un sistema ortonormal, entonces son equivalentes:
    \begin{enumerate}
        \item $S$ es completo.
        \item Si $x \in \mathcal{H}$ y $(x, e_i) = 0 \quad \forall i \in \mathcal{I} \implies x = 0$.
        \item Para cada $x \in \mathcal{H}$ se cumple
        $$x = \sum_{i \in \mathcal{I}} (x, e_i) e_i$$
        \item (\textbf{Identidad de Parseval}) Para cada $x \in \mathcal{H}$
        $$\|x\|^2 = \sum_{i \in \mathcal{I}} |(x, e_i)|^2$$
    \end{enumerate}
\begin{proof}
    \begin{description}
        \item[$1 \Longrightarrow 2)$] Ya visto anteriormente.
        \item[$2 \Longrightarrow 3)$] Sabemos que $\sum_{i \in \mathcal{I}} (x, e_i) e_i$ es convergente en $\mathcal{H}$. Sea $z := \sum_{i \in \mathcal{I}} (x, e_i) e_i$. Mostremos que $z = x$:
        Sea $S(x) = \{ i \in \mathcal{I} \mid (x, e_i) \neq 0 \}$. Entonces:
        \[ (x - z, e_i) = \lim_{m \to +\infty} (x - s_m, e_i) = \lim_{m \to +\infty} [ (x, e_i) - (s_m, e_i) ] \]
        Como $(s_m, e_i) = \left( \sum_{j=1}^m (x, e_j) e_j, e_i \right) = (x, e_i) \quad \forall m \ge i$:
        \[ \implies (x - z, e_i) = 0 \quad \forall m \ge i, \ \forall i \implies x - z = 0 \implies x = z. \]
        \item[$3 \Longrightarrow 4)$] $x = \sum_{i=1}^{+\infty} (x, e_i) e_i \implies x = \lim_{m \to +\infty} s_m$, entonces:
        \[ \|x\|^2 = \lim_{m \to +\infty} \|s_m\|^2 = \lim_{m \to +\infty} \sum_{i,j=1}^m (x, e_i)\overline{(x, e_j)} (e_i, e_j) = \lim_{m \to +\infty} \sum_{i=1}^m |(x, e_i)|^2 = \sum_{i=1}^{+\infty} |(x, e_i)|^2 \]
        \item[$4 \Longrightarrow 1)$] Supongamos por el contrario (P.A.) que $S$ no es completo $\implies \exists S' \supsetneq S$ y por lo tanto $\exists z \in \mathcal{H} : z \neq 0$ y $(z, e_i) = 0 \quad \forall i \in \mathcal{I}$. Entonces:
        \[ \|z\|^2 = \sum_{i \in \mathcal{I}} \underbrace{|(z, e_i)|^2}_{0} = 0 \quad \text{pero } z \neq 0 \quad \textbf{ABSURDO}. \]
    \end{description}
\end{proof}
\end{prop}

\begin{teo} 
    Sea $\mathcal{H}$ un espacio de Hilbert, entonces
    \begin{equation*}
        \mathcal{H} \text{ es separable} \iff \text{ existe un sistema ortonormal } S \text{ completo y numerable}
    \end{equation*}
\begin{proof} % //TODO: leer demostración
    \begin{description}
        \item[$\Longleftarrow)$] Si existe $S$ numerable y completo $\implies S = \{e_1, \dots, e_m, \dots\}$ es numerable y completo: $\overline{[S]} = \mathcal{H}$.
        Tomamos $\overline{\{ \text{combinaciones lineales de } S \text{ con coeficientes racionales} \}} = \mathcal{H}$.
        \item[$\Longrightarrow)$] Todo sistema $S$ ortonormal es numerable. Si $z, z' \in S : z \neq z'$, entonces:
        \[ \|z - z'\|^2 = \underbrace{\|z\|^2}_{1} + \underbrace{\|z'\|^2}_{1} - \underbrace{2(z, z')}_{0} = 2 \implies \|z - z'\| = \sqrt{2} \quad \forall z, z' \in S \text{ distintos}. \]
        Entonces:
        Las bolas $B_1(z)$ (bolas de radio 1 centradas en $z$) son disjuntas.

        Por hipótesis $\mathcal{H}$ es separable $\implies \{x_m\}$ es una sucesión densa en $\mathcal{H}$ y, por lo tanto, en cada bola $B_1(z)$ existe un elemento de la sucesión: $x_m \in B_1(z)$.
        Como la sucesión $\{x_m\}$ es numerable, entonces $S$ es numerable.
    \end{description}
\end{proof}
\end{teo}

\begin{teo}[Teorema isomorfismo euclídeo]
    Todo espacio de Hilbert $\mathcal{H}$ separable es isomorfo a $\ell_2$ mediante una isometría.
\begin{proof}
    Sea $S = \{e_k\}_{k=1}^{+\infty}$ un sistema ortonormal de $\mathcal{H}$, entonces $\forall x \in \mathcal{H}$, tenemos 
    $$x = \sum_{i=1}^{+\infty} (x, e_i) e_i \quad \text{ con } \alpha_k := (x, e_k)$$
    Definimos el operador lineal
    \Func{T}{\mathcal{H}}{\ell_2}{x}{(\alpha_1, \alpha_2, \dots)}
    veámos que $T$ es una isometría y, por lo tanto, un isomorfismo. \\
    \begin{itemize}
        \item Veámos que \underline{está bien definido}, entonces tenienod lo siguiente
        \begin{equation*}
            \|(\alpha_1, \alpha_2, \dots)\|_{\ell_2}^2 = \sum_{i=1}^{+\infty} |\alpha_i|^2 = \sum_{i=1}^{+\infty} |(x, e_i)|^2 \le \|x\|^2 < +\infty
        \end{equation*}
        vemos que está bien definido.
        \item Demostremos que \underline{$T$ es una isometría}, entonces
        \begin{equation*}
            \|T(x)\|_{\ell_2} = \|\alpha\|_{\ell_2} = \sqrt{\sum_{i=1}^{+\infty} |\alpha_i|^2} = \sqrt{\sum_{i=1}^{+\infty} |(x, e_i)|^2} = \|x\|
        \end{equation*}
        \item Demostremos que \underline{$T$ es inyectiva}, entonces
        \begin{equation*}
            \|Tx\| = 0 \implies \|\alpha\|_{\ell_2} = 0 \implies \sum_{i=1}^{+\infty} |\alpha_i|^2 = 0 \implies (x, e_i) = 0 \quad \forall i \implies x = 0
        \end{equation*}
        donde queda demostrada la inyectividad. Notar que está manera de demostrar la inyectividad se da gracias a que $T$ es lineal.
        \item Demostremos que \underline{$T$ es sobreyectiva}, por lo que sea $\alpha = (\alpha_1, \alpha_2, \dots) \in \ell_2$, definimos 
        $$x = \sum_{i=1}^{+\infty} \alpha_i e_i$$ 
        Si $x \in \mathcal{H}$, hemos terminado, veámos para ello las sumas parciales
        \begin{align*}
            S_m := \sum_{i=1}^{m} \alpha_i e_i \implies \|S_m - S_n\|^2 &= \left\| \sum_{i=n+1}^{m} \alpha_i e_i \right\|^2 = \sum_{i=n+1}^{m} \sum_{k=n+1}^{m} (\alpha_i e_i, \alpha_k e_k)= \\
            &= \sum_{i=n+1}^{m} \sum_{k=n+1}^{m} \alpha_i \bar{\alpha}_k (e_i, e_k)= \sum_{i=n+1}^{m} |\alpha_i|^2 \xrightarrow{m, n \to +\infty} 0
        \end{align*}
    \end{itemize}
\end{proof}
\end{teo}

\begin{ejemplo}
    $L^2((0, 2\pi))$ es un espacio de Hilbert separable, entonces es isométrico a $\ell_2$. \\ \\
    \noindent
    $\alpha = (\alpha_1, \alpha_2, \dots) \in \ell_2 \implies \exists g \in L^2((0, 2\pi))$ tal que $\alpha_j := (g, e_j)$ para alguna base $\{e_j\}$ de $L^2$ (por ejemplo, funciones trigonométricas).
\end{ejemplo}
